
% Document settings
\documentclass[11pt]{article}
\usepackage[margin=1in]{geometry}
\usepackage{graphicx}
\usepackage{wrapfig}
\usepackage{multirow}
\usepackage{setspace}
\usepackage{hyperref}
\usepackage{xurl}
\usepackage{booktabs}
\pagestyle{plain}
\setlength\parindent{0pt}

\begin{document}

\Huge INF632 Final Presentation
\normalsize ~

\noindent\rule{\textwidth}{1pt}\vspace{.5em}

\LARGE Preparing a Slide Deck to Support Your Presentation
\normalsize ~

The crowning effort of your project this semester will be a 10 minute presentation on your project.  Here is an outline to get you started, but remember, you need not follow this exactly.  What you need to do, is tell your story.

\begin{enumerate}
	\item Opening — about one minute
	\begin{itemize}
		\item {\em Hook:} one clear sentence about the problem or need your wearable addresses.
		\item {\em Takeaway:} one sentence stating your main result or contribution.
	\end{itemize}
	\item Motivation \& Background — about a minute and a half
	\begin{itemize}
		\item {\em Why it matters:} brief context, target users, and why existing solutions fall short, OR why you had interest in this problem.
		\item {\em Key prior work / concepts:} 1–2 short bullets of what we must know.
	\end{itemize}
	\item Design \& System Overview — maybe two minutes
	\begin{itemize}
		\item {\em High‑level architecture diagram:} hardware + software blocks (sensor(s), micro-controller, communication, power, data pipeline).
		\item {\em Design choices:} 3 concise reasons for the selected sensors, form factor, or algorithms.
	\end{itemize}
	\item Implementation — about a minute and a half
	\begin{itemize}
		\item {\em What you built:} short bullets of electronics, your software, data collection setup.
		\item {\em Critical engineering challenges \& fixes:} 2–3 examples such as, noise, latency, calibration.
	\end{itemize}
	\item Data \& Analysis Methods — about a minute and a half
	\begin{itemize}
		\item {\em Dataset:} how much data, who/what, labeling procedure, collection conditions.
		\item {\em Processing \& algorithms:} workflow steps (filtering, features, model type), evaluation metrics.
	\end{itemize}
	\item Results — about two minutes
	\begin{itemize}
		\item {\em Key quantitative results:} 2–4 metrics or plots (accuracy, latency, battery life) with one-sentence interpretation each.
		\item {\em Representative qualitative result:} one short example or demo note (if time/possible).
	\end{itemize}
	\item Limitations \& Future Work — less than one minute
	\begin{itemize}
		\item {\em Top 3 limitations:} for example, sample size, generalizability, comfort.
		\item {\em Next steps:} such as immediate experiments or improvements.
	\end{itemize}
	\item Conclusion \& Takeaway — less than one minute
	\begin{itemize}
		\item {\em One-line summary:} restate your contribution and the impact of your work.
		\item {\em If time allows:} one-sentence invitation for questions.
	\end{itemize}	
\end{enumerate}

Here's some tips useful for any presentation:
\begin{itemize}
	\item Only provide one idea per slide; use large fonts and clear visuals.
	\item Lead with the takeaway and repeat it at the end.
	\item Practice enough that you can stick to the 10 minutes you are allotted.
	\item Use consistent labels and units on plots. Show uncertainty or error bars when relevant.
	\item If you plan to provide a live demo, have a short recorded video as backup.
\end{itemize}
\end{document}
